\chapter{统一的 token-level Mirror Descent 与 OPD 实例化}
\label{ch:four-methods}

本章只使用一套序列建模符号。给定 prompt $x$，student 生成 rollout $y_{1:H}$；第 $t$ 个 token 的历史为 $s_t=(x,y_{<t})$。固定第 $k$ 个 outer step 的 anchor policy $\pi_k$，候选 student policy 记为 $\pi_\theta$，teacher policy 记为 $T$，自然梯度步长统一记为 $\eta>0$。$s_t$ 只表示 rollout 的条件历史，训练样本始终是一条完整的 prompt--rollout；四种 inner 的区别只来自求解同一个 mirror 子问题的方式。

\section{共同的 mirror 子问题与精确解}
\label{sec:ch2-common}

\begin{definition}[Token-level mirror 子问题]
\label{def:ch2-mirror-subproblem}
给定固定历史 $s$、token 信号 $f(s,a)$ 和 anchor $\pi_k$，定义
\begin{equation}
\label{eq:ch2-mirror}
\mathcal M_s(p;f,\pi_k)
=\mathbb E_{a\sim p(\cdot\mid s)}[f(s,a)]
-\frac1\eta\operatorname{KL}\!\left(p(\cdot\mid s)\middle\Vert\pi_k(\cdot\mid s)\right).
\end{equation}
镜像子问题是在所有满足 $p\ll\pi_k$ 的概率分布上最大化 $\mathcal M_s$。同时定义配分函数
\begin{equation}
\label{eq:ch2-mirror-normalizer}
Z_f(s)=\sum_{v\in\mathcal V}\pi_k(v\mid s)\exp\!\bigl(\eta f(s,v)\bigr).
\end{equation}
\end{definition}

\begin{theorem}[KL mirror 子问题的 Gibbs 解]
\label{thm:ch2-gibbs}
在有限词表、$Z_f(s)<\infty$ 且 $\pi_k(v\mid s)>0$ 的条件下，式 \eqref{eq:ch2-mirror} 的唯一解是
\begin{equation}
\label{eq:ch2-optimal-policy}
\pi_k^+(v\mid s)=\frac{\pi_k(v\mid s)e^{\eta f(s,v)}}{Z_f(s)},
\qquad
\log\frac{\pi_k^+(v\mid s)}{\pi_k(v\mid s)}
=\eta f(s,v)-\log Z_f(s).
\end{equation}
对任意候选 $p$，还有精确分解
\begin{equation}
\label{eq:ch2-kl-decomp}
\mathcal M_s(p;f,\pi_k)
=\frac1\eta\log Z_f(s)
-\frac1\eta\operatorname{KL}\!\left(p\middle\Vert\pi_k^+(\cdot\mid s)\right).
\end{equation}
\end{theorem}

\begin{proof}
由式 \eqref{eq:ch2-optimal-policy}，有
\[
\log\frac{p(v\mid s)}{\pi_k^+(v\mid s)}
=\log\frac{p(v\mid s)}{\pi_k(v\mid s)}-\eta f(s,v)+\log Z_f(s).
\]
对 $v$ 按 $p$ 求期望，得到
\[
\operatorname{KL}(p\Vert\pi_k^+)
=\operatorname{KL}(p\Vert\pi_k)-\eta\mathbb E_p[f]+\log Z_f(s).
\]
移项即得式 \eqref{eq:ch2-kl-decomp}。KL 非负且仅在 $p=\pi_k^+$ 时取零，因此最优解唯一。
\end{proof}

\begin{definition}[中心化 token 信号]
\label{def:ch2-centered-signal}
对任意 $f$ 定义 anchor 条件均值和中心化信号
\[
\bar f_k(s)=\mathbb E_{a\sim\pi_k}[f(s,a)\mid s],
\qquad
\widetilde f_k(s,a)=f(s,a)-\bar f_k(s).
\]
若 $f=A^{\pi_k}$ 是 anchor 的优势函数，则 $\bar f_k(s)=0$ 是优势函数定义推出的恒等式；对一般 teacher score 或 OPD log-ratio，$\bar f_k(s)=0$ 不是自动成立的假设。
\end{definition}

\begin{proposition}[优势函数的 anchor 条件均值]
\label{prop:ch2-advantage-centering}
若 $f(s,a)=A^{\pi_k}(s,a)$，则对每个满足 $\pi_k(\cdot\mid s)$ 定义的历史，
\begin{equation}
\label{eq:ch2-advantage-zero-mean}
\mathbb E_{a\sim\pi_k(\cdot\mid s)}[f(s,a)\mid s]=0.
\end{equation}
对于一般 OPD log-ratio 或 teacher score，该等式不是定义恒等式，必须显式估计并减去条件均值。
\end{proposition}

\begin{proof}
由 $A^{\pi_k}(s,a)=Q^{\pi_k}(s,a)-V^{\pi_k}(s)$ 以及
$V^{\pi_k}(s)=\mathbb E_{a\sim\pi_k}[Q^{\pi_k}(s,a)\mid s]$，有
\begin{align*}
\mathbb E_{\pi_k}[A^{\pi_k}(s,a)\mid s]
&=\mathbb E_{\pi_k}[Q^{\pi_k}(s,a)\mid s]-V^{\pi_k}(s)\\
&=V^{\pi_k}(s)-V^{\pi_k}(s)=0.
\end{align*}
\end{proof}

\begin{proposition}[中心化不改变 mirror 最优策略]
\label{prop:ch2-centering}
用 $\widetilde f_k$ 替代 $f$ 不改变式 \eqref{eq:ch2-mirror} 的最优分布。
\end{proposition}

\begin{proof}
因为 $\bar f_k(s)$ 与 action 无关，
\[
\pi_k(v\mid s)e^{\eta\widetilde f_k(s,v)}
=e^{-\eta\bar f_k(s)}\pi_k(v\mid s)e^{\eta f(s,v)}.
\]
分子和归一化常数同时乘以同一个因子 $e^{-\eta\bar f_k(s)}$，归一化后的分布不变。
\end{proof}

\section{从镜像最优性方程到四种可训练目标}
\label{sec:ch2-objective-derivation}

定理 \ref{thm:ch2-gibbs} 给出的不是一个抽象的分布结论，而是一个可以被监督学习近似的逐 token 方程。把式 \eqref{eq:ch2-optimal-policy} 移项，得到
\begin{equation}
\label{eq:ch2-logratio-equation}
\log\pi_k^+(a\mid s)-\log\pi_k(a\mid s)
=\eta f(s,a)-\log Z_f(s).
\end{equation}
左侧是 candidate student 相对于冻结 anchor 的 log-ratio，右侧只由冻结的 token signal 和归一化常数决定。

\begin{definition}[Mirror log-ratio target]
\label{def:ch2-logratio-target}
定义候选 student 相对于冻结 anchor 的 log-ratio 与 Gibbs target 为
\[
U_\theta(s,a)=\log\pi_\theta(a\mid s)-\log\pi_k(a\mid s),
\qquad
H_k(s,a)=\eta f(s,a)-\log Z_f(s).
\]
式 \eqref{eq:ch2-optimal-policy} 等价于 $U_\theta=H_k$；因此 $H_k$ 是参数化 inner 所要实现的逐 token 最优性方程。
\end{definition}

由这个定义，求解 mirror 子问题在参数化策略类中的直接替代，就是在训练测度下让 $U_\theta$ 拟合 $H_k$：
\begin{equation}
\label{eq:ch2-regression-from-mirror}
\min_\theta\ \mathbb E\left[(U_\theta-H_k)^2\right].
\end{equation}
因此 PMD-Part 的平方损失不是额外假设，而是把镜像最优性方程变成参数化模型可优化的投影问题。对式 \eqref{eq:ch2-regression-from-mirror} 求导，使用
\[
\nabla_\theta U_\theta(s,a)=\nabla_\theta\log\pi_\theta(a\mid s),
\]
立即得到式 \eqref{eq:ch2-part-gradient}。

PMD-Mean 来自同一个方程的中心化。由于 $\log Z_f(s)$ 只依赖 $s$，对 anchor action 取条件均值时
\[
\mathbb E_{\pi_k}[H_k(s,a)\mid s]
=\eta\mathbb E_{\pi_k}[f(s,a)\mid s]-\log Z_f(s).
\]
减去该均值后，未知的 $\log Z_f$ 消失，得到
\begin{equation}
\label{eq:ch2-centered-logratio-equation}
H_k(s,a)-\mathbb E_{\pi_k}[H_k(s,\cdot)\mid s]
=\eta\widetilde f_k(s,a).
\end{equation}
这解释了 PMD-Mean 的 target 来源。它不能只回归中心化 target，还必须保留 $U_\theta=\log(\pi_\theta/\pi_k)$ 的概率归一化约束；否则相差一个依赖历史的常数的所有 log-ratio 都会被视为同一个解。

BPO 再对式 \eqref{eq:ch2-logratio-equation} 两边作同样的中心化，但中心化的是 residual：
\begin{align}
D_\theta(s,a)
&=\bigl[\eta f(s,a)-U_\theta(s,a)\bigr]
-\mathbb E_{\pi_k}\bigl[\eta f(s,\cdot)-U_\theta(s,\cdot)\mid s\bigr].
\label{eq:ch2-bpo-from-optimality}
\end{align}
当 $D_\theta=0$ 时，只能先推出 $U_\theta=\eta f-C(s)$；概率归一化再唯一确定 $C(s)=\log Z_f(s)$。这就是 BPO 与 PMD-Part 共享精确零点、但有限参数更新不同的原因。BPO 的 score centering 并非预先加入的 baseline，而是对 $\mathbb E_{\pi_k}[U_\theta]$ 求导后必然产生的项。

MDPO 则不把式 \eqref{eq:ch2-logratio-equation} 改写成回归，而直接对原始 mirror objective 求最大化。由式 \eqref{eq:ch2-kl-decomp}，对任意参数化 student 有
\begin{equation}
\label{eq:ch2-mdpo-kl-gap}
\mathcal J_{\mathrm{MDPO}}(\theta)
=\frac1\eta\log Z_f(s)
-\frac1\eta\operatorname{KL}\!\left(\pi_\theta(\cdot\mid s)\middle\Vert\pi_k^+(\cdot\mid s)\right).
\end{equation}
所以 MDPO 的 inner 是在直接最小化 student 与 Gibbs target 的 reverse KL gap；PMD-Part、PMD-Mean 和 BPO 则是在不同坐标和不同投影测度下最小化该最优性方程的平方残差。

\begin{proposition}[平方回归的投影含义]
\label{prop:ch2-regression-projection}
若 student 函数类不能表示 $H_k$，PMD-Part 的精确 inner 解是 $H_k$ 在函数类诱导的 $L_2(d_{\pi_k}\pi_k)$ 几何下的投影；若训练样本的历史访问测度改为 $m$，投影也相应变为 $L_2(m\pi_k)$。因此 inner loss 为零只在 target 可表示且训练测度覆盖目标测度时才意味着精确 mirror 解。
\end{proposition}

\begin{proof}
平方损失是 Hilbert 空间中的正交投影问题。对任意可行扰动 $v$，最优点 $U^\star$ 满足
\[
\mathbb E_{m,\pi_k}[(U^\star-H_k)v]=0.
\]
这正是 $H_k$ 到可行函数集合的正交投影条件。若 $H_k$ 属于该集合，投影残差为零；否则残差由函数类和测度共同决定。
\end{proof}

\section{四种 inner 的统一推导}
\label{sec:ch2-four-methods}

下面固定一个历史，令
\[
\ell^T(a)=\log T(a\mid s),\qquad
\ell^k(a)=\log\pi_k(a\mid s),\qquad
\ell^\theta(a)=\log\pi_\theta(a\mid s).
\]
训练时只对 $\ell^\theta$ 反向传播；teacher、anchor、目标常数和所有估计的 KL 都停止梯度。

\subsection{PMD-Part：回归完整 Gibbs log-ratio}
\label{sec:ch2-pmd-part}

定义显式 target
\[
H_k(s,a)=\eta f(s,a)-\log Z_f(s),
\qquad
U_\theta(s,a)=\log\pi_\theta(a\mid s)-\log\pi_k(a\mid s).
\]
PMD-Part 在 anchor 测度下最小化
\begin{equation}
\label{eq:ch2-part-loss}
\mathcal L_{\mathrm{Part}}(\theta)
=\frac12\mathbb E_{s\sim d_k,\,a\sim\pi_k}
\left[U_\theta(s,a)-H_k(s,a)\right]^2.
\end{equation}
对 target 停止梯度，链式法则给出
\begin{equation}
\label{eq:ch2-part-gradient}
-\nabla_\theta\mathcal L_{\mathrm{Part}}
=\mathbb E_{d_k,\pi_k}
\left[H_k(s,a)-U_\theta(s,a)\right]
\nabla_\theta\log\pi_\theta(a\mid s).
\end{equation}
这一步是普通的 log-ratio regression；它没有把 KL 视为常数，也没有使用 action-level importance ratio。

\subsection{PMD-Mean：中心化 target 与归一化约束}
\label{sec:ch2-pmd-mean}

由定义可得
\[
\mathbb E_{a\sim\pi_k}[f(s,a)]=\bar f_k(s),
\qquad
\mathbb E_{a\sim\pi_k}[\widetilde f_k(s,a)]=0.
\]
PMD-Mean 用 $\eta\widetilde f_k$ 作为回归 target，并要求 student log-ratio 满足归一化约束
\begin{equation}
\label{eq:ch2-mean-loss}
\min_{U}\ \frac12\mathbb E_{\pi_k}\left[(U-\eta\widetilde f_k)^2\mid s\right]
\quad\text{s.t.}\quad
\mathbb E_{\pi_k}[e^{U}\mid s]=1.
\end{equation}

\begin{theorem}[PMD-Mean 的逐历史 Lambert-$W$ 投影]
\label{thm:ch2-mean-lambert}
在有限词表、$\pi_k(a\mid s)>0$ 且约束可行时，上述中心化平方回归的唯一 KKT 解满足
\begin{equation}
\label{eq:ch2-mean-lambert-theorem}
U_a=G_a-W\!\left(\lambda(s)e^{G_a}\right),\qquad
G_a=\eta\widetilde f_k(s,a),
\end{equation}
其中 $\lambda(s)\geq0$ 由 $\sum_a\pi_k(a\mid s)e^{U_a}=1$ 唯一确定。
\end{theorem}

\begin{proof}
以下给出固定历史 $s$ 的 KKT 推导。令词表上的函数值为 $U_a$ 和 $G_a=\eta\widetilde f_k(s,a)$，拉格朗日函数为
\[
\mathscr L(U,\lambda)
=\frac12\sum_a\pi_k(a\mid s)(U_a-G_a)^2
+\lambda\left(\sum_a\pi_k(a\mid s)e^{U_a}-1\right).
\]
一阶条件为
\[
U_a-G_a+\lambda e^{U_a}=0.
\]
令 $V_a=G_a-U_a$，则 $V_a e^{V_a}=\lambda e^{G_a}$，所以
\begin{equation}
\label{eq:ch2-mean-lambert}
U_a=G_a-W\!\left(\lambda e^{G_a}\right),
\end{equation}
其中 $\lambda$ 由归一化约束唯一确定。为说明唯一性，定义
\[
\Phi(\lambda)=\sum_a\pi_k(a\mid s)\exp\!\left(G_a-W(\lambda e^{G_a})\right),qquad \lambda\geq0.
\]
Lambert--$W$ 在非负轴上单调递增，因此每一项随 $\lambda$ 单调递减；当 $\lambda=0$ 时 $\Phi(0)=\sum_a\pi_k(a\mid s)e^{G_a}$，而 $\lambda\to\infty$ 时 $\Phi(\lambda)\to0$。只要中心化 target 不已经满足 $\Phi(0)=1$，方程 $\Phi(\lambda)=1$ 有唯一正根；若 $\Phi(0)=1$，则 $\lambda=0$，得到 $U_a=G_a$。因此 PMD-Mean 是中心化 target 在指数归一化集合上的投影。严格凸性和 $\Phi$ 的严格单调性给出 KKT 解的唯一性。
\end{proof}

约束 $\sum_a\pi_\theta(a\mid s)=1$ 正是 log-ratio 坐标中的指数归一化条件。特别地，精确 Gibbs 解的 log-ratio 为
\[
U_a^\star=\eta f(s,a)-\log Z_f(s),\qquad
U_a^\star-\eta\widetilde f_k(s,a)=\eta\bar f_k(s)-\log Z_f(s),
\]
它与 action 无关，且自动满足 $\sum_a\pi_k(a\mid s)e^{U_a^\star}=1$；这说明 PMD-Mean 恢复的是同一个 Gibbs 分布，而不是把 $\log Z_f$ 任意设为零。


\subsection{BPO：中心化最优性 residual}
\label{sec:ch2-bpo}

BPO 不显式回归 $\log Z_f$，而将同一个最优性方程的 action-independent 部分消去。定义
\begin{equation}
\label{eq:ch2-bpo-residual}
D_\theta(s,a)
=\eta f(s,a)-U_\theta(s,a)
-\mathbb E_{v\sim\pi_k}[\eta f(s,v)-U_\theta(s,v)\mid s].
\end{equation}
直接展开可得
\[
D_\theta(s,a)
=\eta\widetilde f_k(s,a)-U_\theta(s,a)
+\mathbb E_{v\sim\pi_k}[U_\theta(s,v)\mid s].
\]
未线性化 BPO 的损失为
\begin{equation}
\label{eq:ch2-bpo-loss}
\mathcal L_{\mathrm{BPO}}(\theta)
=\frac1{2\eta}\mathbb E_{d_k,\pi_k}[D_\theta(s,a)^2].
\end{equation}
由于
\[
\nabla_\theta\mathbb E_{v\sim\pi_k}[U_\theta(s,v)]
=\mathbb E_{v\sim\pi_k}[\nabla_\theta\log\pi_\theta(v\mid s)],
\]
所以
\begin{equation}
\label{eq:ch2-bpo-gradient}
-\eta\nabla_\theta\mathcal L_{\mathrm{BPO}}
=\mathbb E_{d_k,\pi_k}\left[D_\theta(s,a)
\left(\nabla_\theta\log\pi_\theta(a\mid s)
-\mathbb E_{v\sim\pi_k}[\nabla_\theta\log\pi_\theta(v\mid s)\mid s]\right)\right].
\end{equation}
BPO 的 score centering 来自对当前 student--anchor KL 的精确求导；它不是因为先对 reward 做了一个近似。

\subsection{MDPO：直接优化 mirror objective}
\label{sec:ch2-mdpo-generic}

MDPO 直接最大化式 \eqref{eq:ch2-mirror} 的参数化版本：
\begin{equation}
\label{eq:ch2-mdpo-generic-objective}
\mathcal J_{\mathrm{MDPO}}(\theta)
=\mathbb E_{a\sim\pi_\theta}[f(s,a)]
-\frac1\eta\operatorname{KL}\!\left(\pi_\theta(\cdot\mid s)\middle\Vert\pi_k(\cdot\mid s)\right).
\end{equation}
完整词表时 KL 是
\begin{equation}
\label{eq:ch2-mdpo-full-kl}
\operatorname{KL}(\pi_\theta\Vert\pi_k)
=\sum_a\pi_\theta(a\mid s)
\left[\log\pi_\theta(a\mid s)-\log\pi_k(a\mid s)\right].
\end{equation}
若 action 从 anchor 采样，定义 $R_\theta(s,a)=\exp(\ell^\theta(a)-\ell^k(a))$，则 action-level IS 形式为
\begin{equation}
\label{eq:ch2-mdpo-sample-integrand}
\widehat{\mathcal J}_{\mathrm{MDPO}}
=\mathbb E_{a\sim\pi_k}\left[R_\theta(s,a)f(s,a)-\frac1\eta R_\theta(s,a)\log R_\theta(s,a)\right].
\end{equation}
其 ascent direction 为
\begin{equation}
\label{eq:ch2-mdpo-generic-gradient}
\nabla_\theta\mathcal J_{\mathrm{MDPO}}
=\mathbb E_{\pi_k}\left[R_\theta(s,a)\left(f(s,a)-\frac1\eta\log R_\theta(s,a)\right)\nabla_\theta\log\pi_\theta(a\mid s)\right].
\end{equation}
额外的 $-R_\theta/\eta$ 项乘以 score 后，在完整 action 求和中因 $\mathbb E_{\pi_\theta}[\nabla\log\pi_\theta]=0$ 消失。MDPO 因而直接优化 mirror objective，其他三种方法优化的是由同一最优性方程构造的 regression objective。

\begin{theorem}[四种 inner 的共同精确解]
\label{thm:ch2-inner-common-root}
在 tabular policy、完整 action 支持和精确词表期望下，PMD-Part、PMD-Mean、未线性化 BPO 和 MDPO 都以式 \eqref{eq:ch2-optimal-policy} 为精确解。PMD-Mean 的结论要求其投影 target 可实现；参数化 student 或有限样本时，四者得到的是不同的投影或优化误差。
\end{theorem}

\begin{proof}
PMD-Part 在零损失处满足 $U_\theta=\eta f-\log Z_f$，于是得到式 \eqref{eq:ch2-optimal-policy}。BPO 的零 residual 只确定 $U_\theta=\eta f-C(s)$；代入 $\mathbb E_{\pi_k}[e^{U_\theta}]=1$ 得 $C(s)=\log Z_f(s)$。PMD-Mean 若中心化 target 属于归一化约束集合，则其平方投影残差为零，并由中心化不变性得到同一策略。MDPO 由式 \eqref{eq:ch2-kl-decomp} 直接得到唯一最大点。
\end{proof}


\begin{proposition}[未线性化 BPO 与 PMD-Part 的梯度差异]
\label{prop:ch2-bpo-part-gradient-gap}
固定一个历史 $s$，令
\begin{equation}
\label{eq:ch2-bpo-part-qc}
q(s,a)=\eta f(s,a)-\log Z_f(s)-U_\theta(s,a),
\qquad
c(s)=\log Z_f(s)+\mathbb E_{\pi_k}[U_\theta(s,\cdot)\mid s]-\eta\mathbb E_{\pi_k}[f(s,\cdot)\mid s].
\end{equation}
则未线性化 BPO residual 满足 $D_\theta=q+c$，并且
\begin{equation}
\label{eq:ch2-bpo-part-gradient-gap}
-\eta\nabla_\theta\mathcal L_{\mathrm{BPO}}
=-\nabla_\theta\mathcal L_{\mathrm{Part}}
+c(s)\,\mathbb E_{a\sim\pi_k}[\nabla_\theta\log\pi_\theta(a\mid s)\mid s].
\end{equation}
因此两者在精确镜像解处都有零梯度，但在一般参数点不相同。它们只有在
$c(s)\,\mathbb E_{\pi_k}[\nabla_\theta\log\pi_\theta]=0$
时才给出相同的方向。
\end{proposition}

\begin{proof}
从 PMD-Part 的最优性残差出发，先把 BPO 的中心化式逐项展开：
\begin{align*}
D_\theta(s,a)
&=\eta f(s,a)-U_\theta(s,a)
 -\mathbb E_{\pi_k}[\eta f(s,\cdot)-U_\theta(s,\cdot)\mid s]\\
&=\eta f(s,a)-\log Z_f(s)-U_\theta(s,a)
 +\log Z_f(s)+\mathbb E_{\pi_k}[U_\theta(s,\cdot)\mid s]\\
&=q(s,a)+c(s).
\end{align*}
由定义，$\mathbb E_{\pi_k}[D_\theta\mid s]=0$。BPO 的精确梯度为
\begin{align*}
-\eta\nabla_\theta\mathcal L_{\mathrm{BPO}}
&=\mathbb E_{\pi_k}\left[D_\theta
 \left(g_\theta-\bar g_\theta\right)\mid s\right],\\
 g_\theta(a)&=\nabla_\theta\log\pi_\theta(a\mid s),
 \qquad \bar g_\theta=\mathbb E_{\pi_k}[g_\theta\mid s].
\end{align*}
代入 $D_\theta=q+c$，并使用 $\mathbb E_{\pi_k}[D_\theta\mid s]=0$，可得
\begin{align*}
\mathbb E_{\pi_k}[D_\theta(g_\theta-\bar g_\theta)\mid s]
&=\mathbb E_{\pi_k}[(q+c)g_\theta\mid s]
 -\mathbb E_{\pi_k}[D_\theta\mid s]\bar g_\theta\\
&=\mathbb E_{\pi_k}[qg_\theta\mid s]+c(s)\bar g_\theta.
\end{align*}
而 PMD-Part 的梯度式为
$-\nabla_\theta\mathcal L_{\mathrm{Part}}=\mathbb E_{\pi_k}[qg_\theta\mid s]$，故式
\eqref{eq:ch2-bpo-part-gradient-gap}
成立。若 $q=0$，则由归一化约束有 $c=0$，两边同时为零；离开该点时 $c$ 通常非零，差异项保留。
\end{proof}

这个命题说明“都在 token level”只保证它们面对同一个局部镜像根，不保证优化轨迹相同。PMD-Part 回归完整的 action-dependent log-ratio，并把 $\log Z_f$ 当作固定 target；BPO 先消去 action-independent 常数，再对当前 student 的归一化项求导，因此多出 centered score。若使用的参数化使
$\mathbb E_{\pi_k}[\nabla_\theta\log\pi_\theta]=0$
（例如采样测度与当前 student 一致，或在特定 tabular 点上），差异项才会消失。

\section{OPD：从 teacher log-ratio 到 outer--inner 算法}
\label{sec:ch2-opd-instantiation}

OPD 的一个训练样本先采样 prompt $x\sim\mathcal D_X$，再从第 $k$ 个 anchor policy 采样一整条 response rollout
$y_{1:H}\sim\pi_k(\cdot\mid x)$。因此样本的联合概率是
\[
 p_k(x,y_{1:H})=\mathcal D_X(x)\prod_{t=1}^{H}
 \pi_k(y_t\mid x,y_{<t}).
\]
这里的 $y_{<t}$ 只出现在条件概率中；优化和统计的基本单位是一条完整的 prompt--rollout，prompt token 只提供条件，不参加 response loss。teacher 和 student 在这条相同 rollout 上逐 token forward，得到
$T(\cdot\mid x,y_{<t})$、$\pi_k(\cdot\mid x,y_{<t})$ 和候选
$\pi_\theta(\cdot\mid x,y_{<t})$。

OPD 不引入额外的外部 reward，而在第 $k$ 个 outer step 使用 teacher 与当前 anchor 的局部 log-ratio：
\begin{equation}
\label{eq:ch2-opd-current-reward}
r_k(s,a)=\log T(a\mid s)-\log\pi_k(a\mid s).
\end{equation}
它是序列 log-ratio 的逐 token 分解：
\begin{equation}
\label{eq:ch2-opd-sequence-telescope}
\sum_{t=1}^{H}r_k(s_t,a_t)
=\log T(y\mid x)-\log\pi_k(y\mid x).
\end{equation}
因此 $r_k$ 是对 teacher--student 序列差异的无损分摊，不是一个与历史无关的固定终端 reward。

\begin{definition}[OPD mirror target]
\label{def:ch2-opd-target}
将 $f=r_k$ 代入式 \eqref{eq:ch2-optimal-policy}，定义
\begin{equation}
\label{eq:ch2-opd-geometric-bridge}
\pi_{k+1}^{\star}(a\mid s)
=\frac{\pi_k(a\mid s)^{1-\eta}T(a\mid s)^\eta}{Z_k(s)},
\qquad
Z_k(s)=\sum_v\pi_k(v\mid s)^{1-\eta}T(v\mid s)^\eta.
\end{equation}
\end{definition}

\begin{theorem}[OPD mirror step 的几何收敛]
\label{thm:ch2-opd-geometric-bridge}
若 $0<\eta\leq1$，则精确 outer 更新满足
\begin{equation}
\label{eq:ch2-opd-projective-contraction}
\log\frac{\pi_{k+1}^{\star}(a\mid s)}{T(a\mid s)}
=(1-\eta)\log\frac{\pi_k(a\mid s)}{T(a\mid s)}-\log Z_k(s).
\end{equation}
因此任意两个 action 的 log-ratio 差满足
\[
\max_a\log\frac{\pi_{k+1}^{\star}(a\mid s)}{T(a\mid s)}
-\min_a\log\frac{\pi_{k+1}^{\star}(a\mid s)}{T(a\mid s)}
=(1-\eta)\Delta_k(s).
\]
$\eta=1$ 时一步得到 $T$；$0<\eta<1$ 时 projective 误差按 $1-\eta$ 几何衰减。
\end{theorem}

\begin{proof}
由式 \eqref{eq:ch2-opd-geometric-bridge}，
\[
\log\frac{\pi_{k+1}^{\star}(a\mid s)}{T(a\mid s)}
=(1-\eta)\log\frac{\pi_k(a\mid s)}{T(a\mid s)}-\log Z_k(s).
\]
减去任意两个 action 的等式，归一化常数抵消，得到 projective 递推。
\end{proof}

\subsection{prompt--rollout 上的四种显式估计量}
\label{subsec:ch2-opd-explicit-inner}

前面的镜像方程还需要翻译成实际训练代码能够读取的量。第 $k$ 个 outer step 固定 prompt batch，并从 anchor policy $\pi_k$ 为每个 prompt 采样一条或多条 response rollout。对第 $i$ 条 rollout 的第 $t$ 个有效 response token，记
\[
 s_{it}=(x_i,y_{i,<t}),
 \qquad a_{it}=y_{i,t},
\]
并在同一历史上做三次 forward：
\begin{equation}
\label{eq:ch2-opd-three-logprobs}
 \ell^T_{it}=\log T(a_{it}\mid s_{it}),\qquad
 \ell^k_{it}=\log\pi_k(a_{it}\mid s_{it}),\qquad
 \ell^\theta_{it}=\log\pi_\theta(a_{it}\mid s_{it}).
\end{equation}
令 $\mathcal I$ 为所有 response-token mask 的索引集合，$N=|\mathcal I|$。prompt token 只提供条件，不进入下式的 loss 求和。

如果能够访问完整 vocabulary logits，则第 $k$ 个 anchor 与 teacher 的条件 KL 可以直接计算为
\begin{equation}
\label{eq:ch2-opd-full-kls}
 D^{kT}_{it}=\sum_{v\in\mathcal V}\pi_k(v\mid s_{it})
 \log\frac{\pi_k(v\mid s_{it})}{T(v\mid s_{it})},
 \qquad
 D^{k\theta}_{it}=\sum_{v\in\mathcal V}\pi_k(v\mid s_{it})
 \log\frac{\pi_k(v\mid s_{it})}{\pi_\theta(v\mid s_{it})}.
\end{equation}
前者给出 $\pi_k$ 测度下的 teacher log-ratio 均值，后者给出 student log-ratio 的归一化中心。若不能计算完整词表和，只能对同一 prompt 的多条 rollout 采用 Monte Carlo 或 Top-$K$ 尾部估计；中心估计必须停止梯度。

\paragraph*{PMD-Part。}
OPD 的局部 reward 是
\[
 r^k_{it}=\ell^T_{it}-\ell^k_{it},
\]
而镜像子问题的精确配分函数为
\[
 Z_k(s_{it})=\sum_{v\in\mathcal V}
 \pi_k(v\mid s_{it})^{1-\eta}T(v\mid s_{it})^\eta.
\]
因此 PMD-Part 的 sampled regression loss 和 ascent direction 分别为
\begin{equation}
\label{eq:ch2-opd-part-explicit-loss}
 \widehat{\mathcal L}_{k,\mathrm{Part}}
 =\frac1{2N}\sum_{(i,t)\in\mathcal I}
 \left[\ell^\theta_{it}-\ell^k_{it}
 -\eta(\ell^T_{it}-\ell^k_{it})
 +\log Z_k(s_{it})\right]^2,
\end{equation}
\begin{equation}
\label{eq:ch2-opd-part-explicit-gradient}
 -\widehat\nabla_\theta\mathcal L_{k,\mathrm{Part}}
 =\frac1N\sum_{(i,t)\in\mathcal I}
 \left[\eta(\ell^T_{it}-\ell^k_{it})-\ell^\theta_{it}+\ell^k_{it}
 -\log Z_k(s_{it})\right]
 \nabla_\theta\log\pi_\theta(a_{it}\mid s_{it}).
\end{equation}
这里的 $\log Z_k$ 是冻结 target 的一部分。若以 $M$ 个 anchor action 估计
\[
 \widehat Z_{k,M}(s)=\frac1M\sum_{j=1}^M
 \exp\!\left(\eta[\log T(a_j\mid s)-\log\pi_k(a_j\mid s)]\right),
 \qquad a_j\sim\pi_k(\cdot\mid s),
\]
则回归 residual 还会叠加 plug-in 的 $\log\widehat Z_{k,M}-\log Z_k$ 偏差；这部分偏差不能通过增大 inner optimizer steps 消除。

\paragraph*{PMD-Mean。}
由式 \eqref{eq:ch2-opd-full-kls}，有
\[
 \mathbb E_{v\sim\pi_k(\cdot\mid s)}
 [\log T(v\mid s)-\log\pi_k(v\mid s)]
 =-D^{kT}(s).
\]
所以中心化的 OPD reward 可以写为
\[
 \widetilde r^k_{it}=\ell^T_{it}-\ell^k_{it}+D^{kT}_{it}.
\]
PMD-Mean 的可执行平方回归和梯度为
\begin{equation}
\label{eq:ch2-opd-mean-explicit-loss}
 \widehat{\mathcal L}_{k,\mathrm{Mean}}
 =\frac1{2N}\sum_{(i,t)\in\mathcal I}
 \left[\ell^\theta_{it}-\ell^k_{it}
 -\eta(\ell^T_{it}-\ell^k_{it}+D^{kT}_{it})\right]^2,
\end{equation}
\begin{equation}
\label{eq:ch2-opd-mean-explicit-gradient}
 -\widehat\nabla_\theta\mathcal L_{k,\mathrm{Mean}}
 =\frac1N\sum_{(i,t)\in\mathcal I}
 \left[\eta(\ell^T_{it}-\ell^k_{it}+D^{kT}_{it})
 -\ell^\theta_{it}+\ell^k_{it}\right]
 \nabla_\theta\log\pi_\theta(a_{it}\mid s_{it}).
\end{equation}
这个式子把未知的 $\log Z_k$ 换成了 anchor 测度下可估计的中心。它依然隐含归一化约束；如果只回归中心化 reward 而不满足该约束，得到的 log-ratio 只在 action-dependent 部分正确。

\paragraph*{未线性化 BPO。}
BPO 同时中心化 reward residual 和当前 student 的 log-ratio。定义
\begin{equation}
\label{eq:ch2-opd-bpo-explicit-residual}
 \delta^{k,\theta}_{it}
 =\eta(\ell^T_{it}-\ell^k_{it}+D^{kT}_{it})
 -\left(\ell^\theta_{it}-\ell^k_{it}+D^{k\theta}_{it}\right).
\end{equation}
未线性化损失以及对 student 的精确梯度分别为
\begin{equation}
\label{eq:ch2-opd-bpo-explicit-loss}
 \widehat{\mathcal L}_{k,\mathrm{BPO}}
 =\frac1{2\eta N}\sum_{(i,t)\in\mathcal I}
 \left(\delta^{k,\theta}_{it}\right)^2,
\end{equation}
\begin{equation}
\label{eq:ch2-opd-bpo-explicit-gradient}
 -\eta\widehat\nabla_\theta\mathcal L_{k,\mathrm{BPO}}
 =\frac1N\sum_{(i,t)\in\mathcal I}\delta^{k,\theta}_{it}
 \left[
 \nabla_\theta\log\pi_\theta(a_{it}\mid s_{it})
 -\sum_{v\in\mathcal V}\pi_k(v\mid s_{it})
 \nabla_\theta\log\pi_\theta(v\mid s_{it})
 \right].
\end{equation}
第二项正是 reverse-KL 对 student 的导数产生的 centered score。它解释了为什么 BPO 与 PMD-Part 共享精确零点，却在有限 inner 更新中沿不同方向移动。

\paragraph*{MDPO。}
令 action-level ratio 为
\[
 R^\theta_{it}=\exp(\ell^\theta_{it}-\ell^k_{it}).
\]
在 anchor rollout 上，MDPO 的 action-level IS objective 和 ascent estimator 为
\begin{equation}
\label{eq:ch2-opd-mdpo-explicit-objective}
 \widehat{\mathcal J}_{k,\mathrm{MDPO}}
 =\frac1N\sum_{(i,t)\in\mathcal I}R^\theta_{it}
 \left[(\ell^T_{it}-\ell^k_{it})
 -\frac1\eta(\ell^\theta_{it}-\ell^k_{it})\right],
\end{equation}
\begin{equation}
\label{eq:ch2-opd-mdpo-explicit-gradient}
 \widehat\nabla_\theta\mathcal J_{k,\mathrm{MDPO}}
 =\frac1N\sum_{(i,t)\in\mathcal I}R^\theta_{it}
 \left[(\ell^T_{it}-\ell^k_{it})
 -\frac1\eta(\ell^\theta_{it}-\ell^k_{it})\right]
 \nabla_\theta\log\pi_\theta(a_{it}\mid s_{it}).
\end{equation}
这个 ratio 只把 anchor action 测度换成当前 student action 测度。若历史来自 replay，还需单独乘 history/prefix 权重；不能把 action ratio、历史 ratio 和未来 suffix ratio 混成一个权重。

\begin{theorem}[prompt--rollout inner 的共同零点]
\label{thm:ch2-opd-inner-common-root}
在每个有效历史上的词表分布均可表示、$Z_k(s)$ 精确可计算、支持条件成立且优化测度覆盖所有有效 prompt--rollout 历史时，四种 inner 的精确解均满足
\begin{equation}
\label{eq:ch2-opd-inner-common-root}
 \pi_\theta(v\mid x_i,y_{i,<t})
 =\frac{\pi_k(v\mid x_i,y_{i,<t})^{1-\eta}
 T(v\mid x_i,y_{i,<t})^\eta}
 {Z_k(x_i,y_{i,<t})}
\end{equation}
对所有有效 $(i,t)$ 成立。有限 rollout、参数化限制和有限 inner updates 会把逐历史结论变成 rollout-mask 测度下的投影。
\end{theorem}

\begin{proof}
PMD-Part 的平方损失在每个历史达到零时给出
\[
 \ell^\theta(v\mid s)-\ell^k(v\mid s)
 =\eta[\ell^T(v\mid s)-\ell^k(v\mid s)]-\log Z_k(s),
\]
指数化即可得到式 \eqref{eq:ch2-opd-inner-common-root}。PMD-Mean 和 BPO 只消去了 action-independent 的中心项；将它们的零 residual 代入归一化恒等式
\[
 \sum_v\pi_k(v\mid s)\exp(\ell^\theta(v\mid s)-\ell^k(v\mid s))=1
\]
会唯一恢复同一个 $\log Z_k(s)$。MDPO 则由 KL 分解
\[
 \mathcal J_{\mathrm{MDPO}}(\theta)
 =\frac1\eta\log Z_k(s)
 -\frac1\eta\operatorname{KL}(\pi_\theta(\cdot\mid s)\Vert
 \pi_{k+1}^{\star}(\cdot\mid s))
\]
在 KL 唯一取零时得到相同分布。对每个 prompt--rollout 历史逐点应用这些结论即可。
\end{proof}

\subsection{访问测度与有限 inner 误差}
\label{subsec:ch2-opd-access-and-inner-error}

一个 outer step 先冻结 $\pi_k$ 与 $T$，再在固定 target 上运行 $M_k$ 个 inner optimizer updates，最后将得到的 student 记为 $\pi_{k+1}$。$M_k$ 是计算预算，不等于数学上的精确收敛。定义 projective 误差振幅
\[
 \Delta_k(s)=\max_{a}\log\frac{\pi_k(a\mid s)}{T(a\mid s)}
 -\min_{a}\log\frac{\pi_k(a\mid s)}{T(a\mid s)}.
\]

\begin{proposition}[有限 inner 误差的 outer 递推]
\label{prop:ch2-opd-inner-error}
若第 $k$ 个 inner 产生的 student 相对于精确 mirror target 的 projective 振幅误差不超过 $\varepsilon_k(s)$，则
\begin{equation}
\label{eq:ch2-opd-inner-error-recursion}
 \Delta_{k+1}(s)
 \leq (1-\eta)\Delta_k(s)+\varepsilon_k(s).
\end{equation}
因此
\[
 \Delta_K(s)\leq(1-\eta)^K\Delta_0(s)
 +\sum_{j=0}^{K-1}(1-\eta)^{K-1-j}\varepsilon_j(s).
\]
若 $\sup_j\varepsilon_j(s)\leq\bar\varepsilon(s)$ 且 $0<\eta\leq1$，则
\[
 \limsup_{K\to\infty}\Delta_K(s)\leq\frac{\bar\varepsilon(s)}{\eta}.
\]
\end{proposition}

\begin{proof}
精确镜像 target 满足
\[
 \log\frac{\pi_{k+1}^{\star}(a\mid s)}{T(a\mid s)}
 =(1-\eta)\log\frac{\pi_k(a\mid s)}{T(a\mid s)}-\log Z_k(s).
\]
对任意两个 action 作差，$\log Z_k(s)$ 消失，故精确 target 的 projective 振幅是 $(1-\eta)\Delta_k(s)$。有限 inner 误差最多再增加 $\varepsilon_k(s)$，得到式 \eqref{eq:ch2-opd-inner-error-recursion}。对递推式反复代入得到有限 $K$ 展开；当误差有统一上界时使用
\[
 \sum_{j=0}^{K-1}(1-\eta)^j\leq\frac1\eta
\]
即可得到上界。
\end{proof}

若 replay 历史由行为测度 $\nu$ 产生，而目标仍是 anchor 的历史测度 $d_{\pi_k}$，定义
\[
 \Omega_k(s_t)=\frac{d_{\pi_k}(s_t)}{d_\nu(s_t)}.
\]
对于只依赖当前历史和当前 token 的 integrand $Y_t(s_t,a_t)$，有
\begin{equation}
\label{eq:ch2-opd-history-is}
 \mathbb E_{d_{\pi_k},\pi_k}[Y_t]
 =\mathbb E_{d_\nu,\pi_k}[\Omega_k(s_t)Y_t].
\end{equation}

\begin{proposition}[局部 token 目标只需要历史权重]
\label{prop:ch2-local-access-measure}
若 $Y_t$ 只依赖 $(x,y_{<t},y_t)$，且 $d_{\pi_k}\ll d_\nu$，则式 \eqref{eq:ch2-opd-history-is} 是无偏的历史测度修正。它不需要未来 token 的 suffix ratio。
\end{proposition}

\begin{proof}
将目标联合分布按当前历史分解：
\[
 \mathbb E_{d_{\pi_k},\pi_k}[Y_t]
 =\sum_{s_t}d_{\pi_k}(s_t)\sum_a\pi_k(a\mid s_t)Y_t(s_t,a).
\]
代入 $d_{\pi_k}(s_t)=\Omega_k(s_t)d_\nu(s_t)$ 即得式 \eqref{eq:ch2-opd-history-is}。由于 $Y_t$ 不包含未来 token，对未来 token 的边际积分在两边都等于 $1$；再乘 suffix ratio 会改变被估计的 token 目标。
\end{proof}

\section{四种 inner 的可执行估计量、无偏性与方差}
\label{sec:ch2-estimator-theory}

本节把四种 inner 写成可以直接实现的随机估计量。所有公式都在第 $k$ 个 outer step 内理解：$\pi_k$、教师 $T$、token reward $r_k$ 和 replay 行为策略 $b$ 冻结，只有学生参数 $\theta$ 更新。记
\[
 g_\theta(s,a)=\nabla_\theta\log\pi_\theta(a\mid s),
 \qquad
 u_\theta(s,a)=\log\frac{\pi_\theta(a\mid s)}{\pi_k(a\mid s)},
\]
并假设策略支持满足 $\pi_k(a\mid s)>0$ 且 $b(a\mid s)>0$ whenever the corresponding target probability is nonzero。对任意向量估计量 $X$，令
\[
 \operatorname{Var}(X)=\mathbb E\|X-\mathbb E X\|_2^2
 =\mathbb E\|X\|_2^2-\|\mathbb E X\|_2^2.
\]

\subsection{统一的单样本和批量方差公式}
\label{subsec:unified-estimator-variance}

先考虑固定 prefix $s$ 和固定 action 测度 $\nu(\cdot\mid s)$。如果目标梯度可以写成
\[
 G_s=\mathbb E_{a\sim\nu}[Y(s,a)],
\]
用 $B$ 个条件独立样本 $Y_1,\ldots,Y_B$ 的平均值
\[
 \widehat G_s=\frac1B\sum_{i=1}^B Y_i
\]
满足
\begin{equation}
\label{eq:inner-unbiased-mean}
 \mathbb E[\widehat G_s]=G_s,
 \qquad
 \mathbb E\|\widehat G_s-G_s\|_2^2
 =\frac1B\operatorname{Var}(Y).
\end{equation}
这说明增加 batch size 只能降低抽样方差，不能消除由错误访问测度、错误的归一化常数或函数类投影造成的系统偏差。

若样本来自 $b$ 而目标 action 测度是 $\nu$，定义
\[
 \omega_\nu(s,a)=\frac{\nu(a\mid s)}{b(a\mid s)}.
\]
只要 $\mathbb E_b[\omega_\nu\|Y\|_2]<\infty$，则
\[
 G_s=\mathbb E_{a\sim b}[\omega_\nu(s,a)Y(s,a)],
\]
因此 $\omega_\nu Y$ 是无偏估计量。其方差由二阶矩控制：若 $\|Y\|_2\leq C_Y$，则
\begin{equation}
\label{eq:inner-is-second-moment}
 \mathbb E_b\|\omega_\nu Y\|_2^2
 \leq C_Y^2\,\mathbb E_b[\omega_\nu^2]
 =C_Y^2\bigl(1+\chi^2(\nu\Vert b)\bigr).
\end{equation}
因此 action-level IS 的代价由 $1+\chi^2(\nu\Vert b)$ 给出，而不是由 ratio 的均值给出；后者恒等于 $1$，不能反映方差。

当 prefix 也来自 replay 时，目标状态测度为 $d_{\pi_k}$，定义
\[
 W^k_{<t}=\prod_{j<t}\frac{\pi_k(y_j\mid s_j)}{b(y_j\mid s_j)}.
\]
对局部信号 $Y_t$，严格的估计量为 $W^k_{<t}Y_t$。若 $\|Y_t\|_2\leq C_Y$ 且每个 prefix 的条件 action ratio 满足
\[
 \mathbb E_b\left[\left(\frac{\pi_k(a\mid s)}{b(a\mid s)}\right)^2\middle|s\right]
 \leq 1+\chi_{\max}^2,
\]
则递推得到
\begin{equation}
\label{eq:prefix-second-moment}
 \mathbb E_b[(W^k_{<t})^2]
 \leq (1+\chi_{\max}^2)^{t-1},
 \qquad
 \mathbb E_b\|W^k_{<t}Y_t\|_2^2
 \leq C_Y^2(1+\chi_{\max}^2)^{t-1}.
\end{equation}
这就是长 prefix 上方差随时间增长的来源。它与 token-level inner 是否使用 regression 无关，只由访问测度校正决定。

\begin{proposition}[投影误差与抽样方差的正交分解]
\label{prop:projection-sampling-decomposition}
设目标梯度为 $G$，算法实际在函数类 $\mathcal F$ 和采样测度 $m$ 下优化，函数类投影给出的梯度为 $G_{\mathcal F,m}$。令 $\widehat G_B$ 是该投影梯度的无偏 batch 估计量，则
\[
 \mathbb E\|\widehat G_B-G\|_2^2
 =\|G_{\mathcal F,m}-G\|_2^2
 +\frac1B\operatorname{Var}(Y_{\mathcal F,m}).
\]
\end{proposition}

\begin{proof}
写成 $\widehat G_B-G=(\widehat G_B-G_{\mathcal F,m})+(G_{\mathcal F,m}-G)$。第一项的期望为零，故交叉项期望为零；再用式 \eqref{eq:inner-unbiased-mean} 即得结论。第一项是不可由加大 batch 消除的 projection/measure bias，第二项才是 Monte Carlo variance。
\end{proof}

\subsection{PMD-Part：完整 log-ratio target 的回归估计}
\label{subsec:estimator-part}

PMD-Part 的单样本 residual 和梯度分别为
\[
 q_{k,\theta}(s,a)=h_k(s,a)-u_\theta(s,a),
 \qquad
 Y_{\mathrm{Part}}(s,a)=q_{k,\theta}(s,a)g_\theta(s,a),
\]
其中
\[
 h_k(s,a)=\eta r_k(s,a)-\log Z_k(s),
 \qquad
 Z_k(s)=\mathbb E_{a\sim\pi_k}\left[e^{\eta r_k(s,a)}\right].
\]
若 $(s,a)\sim d_{\pi_k}\pi_k$，则
\[
 -\nabla_\theta\mathcal L_{\mathrm{Part}}
 =\mathbb E[q_{k,\theta}g_\theta].
\]
在 replay 上的严格无偏估计量是
\begin{equation}
\label{eq:part-unbiased-estimator}
 \widehat G_{\mathrm{Part}}
 =\frac1B\sum_{i=1}^B
 W^k_{<t_i}\,\frac{\pi_k(a_i\mid s_i)}{b(a_i\mid s_i)}
 q_{k,\theta}(s_i,a_i)g_\theta(s_i,a_i).
\end{equation}
若 action 已按 $\pi_k$ 采样，只保留 prefix weight；若把 replay 测度本身作为训练目标，则去掉 $W^k_{<t_i}$ 得到 replay-projection 梯度。

对于真实 $Z_k$，式 \eqref{eq:part-unbiased-estimator} 在支持和二阶矩条件下无偏，且
\[
 \operatorname{MSE}(\widehat G_{\mathrm{Part}})
 =\frac1B\operatorname{Var}_{d_{\pi_k},\pi_k}(q_{k,\theta}g_\theta)
\]
在 on-policy prefix 下成立。工程上常用 $M$ 个 action 样本估计 $Z_k$：
\[
 \widehat Z_{k,M}(s)=\frac1M\sum_{j=1}^M e^{\eta r_k(s,a_j)},
 \qquad a_j\sim\pi_k(\cdot\mid s).
\]
由于 $\log$ 是凹函数，\[
 \mathbb E[\log\widehat Z_{k,M}(s)]\leq\log Z_k(s),
\]
所以 $-\log\widehat Z_{k,M}$ 对 $-\log Z_k$ 产生向上的有限样本偏差。若 $X=e^{\eta r_k}$ 有有限方差，二阶 delta method 给出
\begin{equation}
\label{eq:logz-delta-method}
 \mathbb E[\log\widehat Z_{k,M}]
 =\log Z_k-\frac{\operatorname{Var}(X)}{2MZ_k^2}+O(M^{-2}),
 \qquad
 \operatorname{Var}(\log\widehat Z_{k,M})
 =\frac{\operatorname{Var}(X)}{MZ_k^2}+O(M^{-2}).
\end{equation}
因此 PMD-Part 的主要额外误差来自 $\log Z$ 的 plug-in，而不是回归梯度本身。

\subsection{PMD-Mean：中心化 target 与归一化约束}
\label{subsec:estimator-mean}

令
\[
 r_{k,0}(s,a)=r_k(s,a)-\bar r_k(s),
 \qquad
 \bar r_k(s)=\mathbb E_{\pi_k}[r_k(s,a)].
\]
PMD-Mean 使用目标 $\eta r_{k,0}$，并在函数空间中施加
\[
 \mathbb E_{\pi_k}[e^{u_\theta(s,a)}\mid s]=1.
\]
对一个 prefix 采 $K$ 个 action，定义
\[
 \widehat{\bar r}_{k,i}=\frac1K\sum_{j=1}^K r_k(s_i,a_{ij}),
 \qquad
 \widehat Y_{ij}^{\mathrm{Mean}}
 =\left[\eta(r_k(s_i,a_{ij})-\widehat{\bar r}_{k,i})-u_\theta(s_i,a_{ij})\right]g_\theta(s_i,a_{ij}).
\]
如果把中心和当前 action 用独立样本计算，则 $\mathbb E[\widehat{\bar r}_{k,i}\mid s_i]=\bar r_k(s_i)$，从而该估计量对中心化 PMD-Mean 梯度无偏。若中心和当前 action 共用同一组 $K$ 样本，标量中心本身仍无偏，但它与 $g_\theta(s_i,a_{ij})$ 相关，导致有限 $K$ 的乘积估计量出现协方差项：
\begin{equation}
\label{eq:mean-covariance-bias}
 \mathbb E\left[(r_j-\widehat{\bar r})g_j\mid s\right]
 =\mathbb E[(r-\bar r)g\mid s]-\frac1K\operatorname{Cov}_{\pi_k}(r,g\mid s).
\end{equation}
一个简单的修正是 leave-one-out center：
\[
 \widehat{\bar r}_{-j}=\frac1{K-1}\sum_{\ell\ne j}r_{\ell},
 \qquad K>1.
\]
此时 $\widehat{\bar r}_{-j}$ 与 $a_j$ 条件独立，故对应的 product estimator 无偏。其方差包含
\[
 \operatorname{Var}(\widehat{\bar r}_{-j}\mid s)
 =\frac{\operatorname{Var}_{\pi_k}(r\mid s)}{K-1},
\]
因此 $K$ 增大可以降低 center noise，但不能消除 model projection bias。

\subsection{BPO：完整中心化 score 的有限样本估计}
\label{subsec:estimator-bpo}

令
\[
 d_{k,\theta}(s,a)=\eta r_k(s,a)-u_\theta(s,a),
 \qquad
 \bar d_{k,\theta}(s)=\mathbb E_{\pi_k}[d_{k,\theta}(s,a)],
\]
则未线性化 BPO 的精确梯度可写为
\[
 G_{\mathrm{BPO}}(s)=
 \mathbb E_{\pi_k}\left[(d_{k,\theta}-\bar d_{k,\theta})
 (g_\theta-\bar g_\theta)\mid s\right],
 \qquad
 \bar g_\theta=\mathbb E_{\pi_k}[g_\theta\mid s].
\]
用同一 prefix 的 $K$ 个 action 直接计算样本中心：
\[
 \widehat d_{ij}=d_{ij}-\frac1K\sum_{\ell=1}^K d_{i\ell},
 \qquad
 \widehat g_{ij}=g_{ij}-\frac1K\sum_{\ell=1}^K g_{i\ell}.
\]
则
\begin{equation}
\label{eq:bpo-sample-covariance}
 \widehat G_{\mathrm{BPO},i}
 =\frac1K\sum_{j=1}^K\widehat d_{ij}\widehat g_{ij}
\end{equation}
是一个样本协方差。条件于 prefix，i.i.d. 样本给出精确恒等式
\begin{equation}
\label{eq:bpo-bias-factor}
 \mathbb E[\widehat G_{\mathrm{BPO},i}\mid s_i]
 =\frac{K-1}{K}G_{\mathrm{BPO}}(s_i).
\end{equation}
因此共用中心时，BPO 的梯度有一个明确的 $-(1/K)$ shrinkage bias；乘以 $K/(K-1)$ 可消除该 bias，但会放大方差。等价地，可以使用一个 action 计算中心、另一个独立 action 计算 residual 和 score。对向量 $\widehat d\widehat g$，若四阶矩有限，其 batch variance 仍按 $1/B$ 缩放；$K$ 只减少每个 prefix 的条件协方差。

\begin{proof}
记 $X_j=d_j-\bar d$、$Y_j=g_j-\bar g$。样本中心化恒等式为
\[
 \frac1K\sum_j(d_j-\bar d_K)(g_j-\bar g_K)
 =\frac1K\sum_jX_jY_j-\left(\frac1K\sum_jX_j\right)\left(\frac1K\sum_jY_j\right).
\]
第一项期望为 $\mathbb E[XY]$，第二项期望为 $K^{-1}\mathbb E[XY]$，因为不同样本的交叉项均为零，故得到式 \eqref{eq:bpo-bias-factor}。
\end{proof}

\subsection{MDPO：action ratio、Top-k KL 与完整梯度}
\label{subsec:estimator-mdpo}

固定 prefix 后，MDPO 的梯度为
\[
 G_{\mathrm{MDPO}}(s)=\mathbb E_{a\sim\pi_\theta}
 \left[c_{k,\theta}(s,a)g_\theta(s,a)\right],
 \quad
 c_{k,\theta}=r_k-\frac1\eta\log\frac{\pi_\theta}{\pi_k}.
\]
若 replay action 来自 $b$，单样本无偏估计量为
\begin{equation}
\label{eq:mdpo-action-is-estimator}
 X_{\mathrm{MDPO}}(s,a)=
 \frac{\pi_\theta(a\mid s)}{b(a\mid s)}c_{k,\theta}(s,a)g_\theta(s,a),
 \qquad a\sim b(\cdot\mid s).
\end{equation}
若状态也需从 $d_b$ 换成 $d_{\pi_k}$，再乘 $W^k_{<t}$。在 $\|c_{k,\theta}g_\theta\|\leq C$ 时，二阶矩满足
\[
 \mathbb E_b\|X_{\mathrm{MDPO}}\|^2
 \leq C^2\bigl(1+\chi^2(\pi_\theta\Vert b)\bigr).
\]

当 token vocabulary logits 可访问时，可以用 Top-k 估计 reverse KL 而不截断尾部。取任意 head 集 $S(s)$（Top-$K$ 只是降低方差的选择），令 $P=\pi_\theta(\cdot\mid s)$、$Q=\pi_k(\cdot\mid s)$，并从 $a\sim P$ 采一个尾部样本。记 $r_\theta(a)=P_a/\operatorname{sg}(P_a)$，前向值为 $1$ 但反向传播保留 $P_a$ 的 score 梯度。定义
\begin{equation}
\label{eq:topk-reverse-kl-estimator}
 \widehat D_{K}(P\Vert Q)
 =\sum_{j\in S}P_j\log\frac{P_j}{Q_j}
 +\mathbf 1\{a\notin S\}\,r(a)\,\operatorname{sg}\!\left(-\log\frac{Q_a}{P_a}\right),
\end{equation}
第一项是 head 的精确 reverse KL；第二项的期望为
\[
 \sum_{a\notin S}P_a\frac{Q_a}{P_a}\log\frac{P_a}{Q_a}
 =\sum_{a\notin S}Q_a\log\frac{P_a}{Q_a},
\]
与第一项相加得到 $\KL(P\Vert Q)$。对 $\theta$ 求梯度时，head 使用精确梯度，tail 使用停止梯度的 ratio 和 score；逐项展开可得其期望等于 $\nabla_\theta\KL(P\Vert Q)$。因此 Top-k 估计量同时满足
\[
 \mathbb E[\widehat D_K]=\KL(P\Vert Q),
 \qquad
 \mathbb E[\nabla_\theta\widehat D_K]=\nabla_\theta\KL(P\Vert Q).
\]
若完全截断 tail，第二个等式右侧少掉 tail 梯度，产生不可由 batch size 消除的 bias。Top-k 的方差仅来自 tail Bernoulli 和 tail action；随着 head mass 增大，tail 二阶矩通常下降。

\begin{proposition}[四种 inner 的无偏条件]
\label{prop:four-inner-unbiased}
在 target、anchor 和 reward 都停止梯度，策略支持覆盖成立，且所有相关一阶和二阶矩有限时：
\begin{enumerate}[label=(\roman*)]
 \item PMD-Part 在使用真实 $Z_k$、正确 prefix/action ratio 时无偏；用 $\widehat Z$ 时产生式 \eqref{eq:logz-delta-method} 的 plug-in bias。
 \item PMD-Mean 在中心估计与当前 action 独立，或使用 leave-one-out center 时无偏；同样本 center 产生式 \eqref{eq:mean-covariance-bias} 的 covariance bias。
 \item BPO 的同 prefix 样本协方差估计量带有 $(K-1)/K$ shrinkage；乘以 $K/(K-1)$ 或使用独立中心后无偏。
 \item MDPO 的 action-level IS 估计量无偏；prefix ratio 用于把 replay occupancy 换成 $d_{\pi_k}$，不能用 full suffix ratio 代替。
\end{enumerate}
\end{proposition}

\section{inner 算法的选择：理论与工程折中}
\label{sec:ch2-inner-choice}

四种方法共享同一个 token-level mirror root，但它们需要的模型访问、随机量和 off-policy 修正不同。对 OPD，教师通常能够在已访问 prefix 上返回 logits，学生也能够返回同一 prefix 的 logits；这个事实决定了默认方案应优先利用 logit 信息，而不是把整个 token 分布重新当作黑盒采样器。

\paragraph{PMD-Part 是默认 inner。}
它直接回归
\[
 u_\theta(s,a)\approx h_k(s,a)=\eta r_k(s,a)-\log Z_k(s),
\]
梯度是标量 residual 与 score 的乘积，不需要 action-level IS（action 按 $\pi_k$ 采样时）。$Z_k$ 可以用全词表 log-sum-exp 精确计算；词表太大时，式 \eqref{eq:topk-reverse-kl-estimator} 提供了 head exact plus tail sample 的无偏替代。它的主要偏差来源只有参数函数类投影和 $\widehat{\log Z}$ 的 plug-in，而不是中心估计与 score 的乘积偏差。因此在 OPD 有 teacher logits、能够缓存 top-$K$ logits 时，PMD-Part 的实现最短，诊断指标也最直接：监控 $\|u_\theta-h_k\|^2$ 和 mirror KL。

\paragraph{MDPO 是需要直接优化时的备选。}
MDPO 不构造 $h_k$，而直接最大化 reward 减 reverse KL。它适合 teacher reward 不便逐 action 回归、但可以稳定计算 candidate policy ratio 的场景。代价是每一个 replay action 都带有 $\pi_\theta/b$，prefix reuse 还要带 $W^k_{<t}$；二阶矩随 action ratio 和 prefix ratio 同时增长。Top-k KL 可以把 reverse-KL 的词表求和代价降到 $O(K)$，但不能消除 action IS 本身的方差。工程上应把 MDPO 放在 PMD-Part 之后作为 ablation 或在 teacher logits 不完整时使用。

\paragraph{PMD-Mean 适合不计算配分函数的实现。}
它只需要估计 $\bar r_k(s)$，但必须维护
\[
 \mathbb E_{\pi_k}[e^{u_\theta}\mid s]=1
\]
这一非线性约束。若使用同一组 action 同时计算中心和 residual，会出现式 \eqref{eq:mean-covariance-bias} 的 covariance bias；leave-one-out 或独立中心可以修复无偏性，却提高每个 prefix 的采样成本。它适合作为不具备完整 vocabulary logits、但能在同一 prefix 采多次 action 的实现。

\paragraph{BPO 适合做低方差中心化对照。}
BPO 的 centered score 去除了 action-independent baseline，理论上能降低梯度方差；但 token-level 实现至少需要每个 prefix 的两个 action，直接使用 $K$ 个样本时还有式 \eqref{eq:bpo-bias-factor} 的 $ (K-1)/K$ shrinkage。BPO 还需要对 residual 和 score 同时做中心化，调试时难以区分 reward noise、中心估计误差和模型投影误差。因而它更适合在 PMD-Part 已经跑通后验证中心化是否带来实际方差收益，而不适合作为第一版 OPD inner。

综合而言，建议的工程顺序是：\textbf{PMD-Part + Top-k KL} 作为主算法，\textbf{MDPO + Top-k reverse KL} 作为直接目标基线，PMD-Mean 和 BPO 作为分别检验“去掉配分函数”和“中心化 score”作用的对照。四者都应使用相同的 prefix replay、相同的 $W^k_{<t}$ 选项和相同的 inner stopping rule，这样差异才归因于 inner 几何而不是访问测度。

\subsection{可执行的 outer--inner 流程}
\label{subsec:outer-inner-algorithm}

第 $k$ 个 outer step 的实现顺序如下。先冻结 $\pi_k$ 和 teacher $T$，在 replay prefix 上计算 $r_k=\log T-\log\pi_k$，并缓存学生、teacher 的 top-$K$ logits。然后选择一种 inner：PMD-Part 回归 $h_k$，PMD-Mean 回归中心化 reward，BPO 回归中心化 residual，MDPO 使用 action-ratio mirror ascent。每次 inner update 都在同一冻结 target 上进行；当估计的 mirror gap、log-ratio residual 或 held-out KL 降到 $\varepsilon_k$ 后，才把所得学生记为 $\pi_{k+1}$ 并刷新 teacher。这样 outer 的 teacher 不会在 inner 尚未稳定时移动，避免把两个时间尺度混成一个 optimizer trajectory。

\section{OPD 的 outer 规划与 inner 容差}
\label{sec:ch2-outer-plan}

一个 outer step 应当把 teacher、anchor 和 inner 优化严格分开。第 $k$ 步先冻结学生 $\pi_k$，由 teacher 产生
\[
 r_k(s,a)=\log\frac{T(a\mid s)}{\pi_k(a\mid s)}.
\]
随后定义镜像目标
\[
 \pi_{k+1}^{\star}(a\mid s)
 =\frac{\pi_k(a\mid s)e^{\eta r_k(s,a)}}{Z_k(s)}
 =\frac{\pi_k(a\mid s)^{1-\eta}T(a\mid s)^\eta}{Z_k(s)}.
\]
在 $0<\eta\leq1$ 时，它是 student 与 teacher 的几何中间分布；$\eta=1$ 才是一步到 teacher。实际算法先固定这个 target，再用四种 inner 之一得到近似参数 $\pi_{k+1}$，最后才进入下一次 teacher refresh。

\subsection{精确 outer 收缩与有限 inner 误差}
\label{subsec:outer-contraction-inner-error}

定义同一 prefix 上任意两个 action 的 projective log-ratio
\[
 \Delta_k(s;a,a')
 =\log\frac{\pi_k(a\mid s)}{T(a\mid s)}
 -\log\frac{\pi_k(a'\mid s)}{T(a'\mid s)}.
\]
精确 mirror update 满足
\begin{equation}
\label{eq:outer-projective-contraction}
 \Delta_{k+1}(s;a,a')=(1-\eta)\Delta_k(s;a,a').
\end{equation}
若 inner 产生的 logit/projection error 为 $e_k(s,a)$，且其 projective oscillation 满足
\[
 \operatorname{osc}_s(e_k)=\max_a e_k(s,a)-\min_a e_k(s,a)\leq\varepsilon_k,
\]
则
\begin{equation}
\label{eq:outer-inexact-recurrence}
 \operatorname{osc}_s(\rho_{k+1})
 \leq (1-\eta)\operatorname{osc}_s(\rho_k)+\varepsilon_k,
 \qquad
 \rho_k=\log(\pi_k/T).
\end{equation}
迭代展开为
\[
 \operatorname{osc}_s(\rho_K)
 \leq (1-\eta)^K\operatorname{osc}_s(\rho_0)
 +\sum_{j=0}^{K-1}(1-\eta)^{K-1-j}\varepsilon_j.
\]
若 $\varepsilon_j\leq\bar\varepsilon$，渐近误差不超过 $\bar\varepsilon/\eta$。因此 outer 目标误差为 $\delta_{\mathrm{out}}$ 时，应把每个 inner 的 projective error 控制在
\[
 \varepsilon_k\lesssim\eta\,\delta_{\mathrm{out}}.
\]
这给出了停止条件，而不是预先规定一个与 $\eta$ 无关的 optimizer step 数。

\subsection{inner 的 Natural Gradient 收敛与步长系数}
\label{subsec:eta-inner-steps}

在第 $k$ 步的局部参数模型中，令 $A_k$ 是 replay prefix 上的 Fisher 曲率，特征值属于 $[m_k,M_k]$。普通梯度对
\[
 Q_k(\Delta)=\frac{1}{2\eta}\Delta^{\mathsf T}A_k\Delta-b_k^{\mathsf T}\Delta
\]
做步长为 $\gamma$ 的更新，其误差传播矩阵为
\[
 E_{k,\eta}=I-\frac{\gamma}{\eta}A_k,
 \qquad
 q_k(\eta,\gamma)=\max\left\{\left|1-\frac{\gamma m_k}{\eta}\right|,
 \left|1-\frac{\gamma M_k}{\eta}\right|\right\}.
\]
收敛要求 $0<\gamma<2\eta/M_k$，并且达到 inner 误差 $\epsilon$ 至少需要
\begin{equation}
\label{eq:inner-step-count}
 N_{\mathrm{in}}
 \geq
 \frac{\log(\epsilon_0/\epsilon)}{-\log q_k(\eta,\gamma)}.
\end{equation}
若把实际 optimizer step 按比例设为 $\gamma=\eta\widetilde\gamma$，则
\[
 q_k=\max\{|1-\widetilde\gamma m_k|,|1-\widetilde\gamma M_k|\},
\]
收敛率与 $\eta$ 一阶解耦；最优固定步长为
\[
 \widetilde\gamma^\star=\frac{2}{m_k+M_k},
 \qquad
 q_k^\star=\frac{M_k-m_k}{M_k+m_k}.
\]
若 $\gamma$ 不随 $\eta$ 缩放，则 $\eta$ 越小，$\gamma/\eta$ 越大，可能越快也可能越过稳定区间；$\eta$ 越大，$\gamma/\eta$ 越小，inner 收敛变慢。这个结论解释了为什么 Natural Gradient 的学习率应与 mirror 系数联动，而不是在所有 outer step 使用固定的 AdamW learning rate。

outer 与 inner 的总成本还受目标移动量影响。对 $0<\eta\ll1$，\[
 \KL(\pi_{k+1}^{\star}\Vert\pi_k)=O(\eta^2),
 \qquad
 \operatorname{osc}(u_k^\star)=O(\eta),
\]
所以小 $\eta$ 每个 outer target 更近，通常需要更少的 inner updates 才达到绝对参数误差；但 outer 收缩率 $1-\eta$ 变慢，需要 $O(\eta^{-1}\log(1/\delta_{\mathrm{out}}))$ 个 outer steps。大 $\eta$ 减少 outer 次数，却放大单次 target 位移、ratio 和 regression residual 的二阶矩；当 $\eta>1$ 时，式 \eqref{eq:outer-projective-contraction} 进入符号翻转区间，$1<\eta<2$ 仍有 projective 绝对收缩但会振荡，$\eta\geq2$ 不再收缩。因此工程默认应取 $0<\eta\leq1$，并用 inner mirror gap 或 KL 停止，而不是固定 inner step 数。

\paragraph{MDPO 的 OPD inner（补充说明）。}
式 \eqref{eq:ch2-opd-mdpo-explicit-objective} 已经给出 action-level IS 版本。若 teacher、anchor 和 student 都能输出完整词表分布，MDPO 也可以直接计算
\[
\sum_{v\in\mathcal V}\pi_\theta(v\mid s_i)
\left[\log T(v\mid s_i)-\log\pi_k(v\mid s_i)\right]
-\frac1\eta\sum_{v\in\mathcal V}\pi_\theta(v\mid s_i)
\log\frac{\pi_\theta(v\mid s_i)}{\pi_k(v\mid s_i)}.
\]
这是一条对每个 prefix 的可微 KL objective；它不需要显式的 sampled action ratio。工程实现通常将第二项写成 logits 上的 reverse-KL penalty，或使用 token-level sampled-KL estimator。只有当数据由旧行为策略 $b$ 产生、而目标 expectation 属于当前 anchor prefix 测度时，才需要额外的 $\omega_i^k$；action-level $\rho_{ij}^{\theta}$ 与 prefix-level $\omega_i^k$ 负责两个不同的测度变换。

\subsection{为什么 inner 需要多步：从精确 KL 几何到 Natural Gradient}
\label{sec:ch2-natural-gradient}

前面的 mirror 解是在策略分布空间中定义的，而实际训练变量是神经网络参数 $\theta$。因此需要把“固定一个 outer target 后的 inner 优化”从分布空间搬到参数空间。下面先不假设局部二次模型，而是从 softmax 参数化和 KL 的定义逐步展开。

固定一个 prefix $s$，令 $z_\theta(s)\in\mathbb R^{|\mathcal A|}$ 为学生 logits，$\pi_\theta(\cdot\mid s)=\operatorname{softmax}(z_\theta(s))$。在第 $k$ 个 outer step 中冻结 $\pi_k$ 和 $r_k$，把 inner 的镜像目标写成最小化形式
\begin{equation}
\label{eq:ch2-opd-inner-objective}
\mathcal J_k(\theta)
=
-\E_{s,a}[r_k(s,a)]
+\frac{1}{\eta}\E_s\!\left[
\KL\!\left(\pi_\theta(\cdot\mid s)\Vert\pi_k(\cdot\mid s)\right)
\right],
\end{equation}
其中 $\E_{s,a}$ 的 action 测度可以是 $\pi_\theta$、$\pi_k$ 或 replay 加权测度；下文的局部曲率只要求该测度固定且具有有限二阶矩。

\paragraph{第一步：logit 空间的 KL 一阶项消失。}
令 $\delta z=z_\theta-z_{\theta_k}$，其中 $\pi_{\theta_k}=\pi_k$。对固定的行为分布 $q=\pi_k$ 定义
\[
K(z)=\KL(\operatorname{softmax}(z)\Vert q).
\]
直接对 softmax 的 log-partition 展开：
\[
K(z_k+\delta z)
=
\sum_a p_a\log p_a-\sum_a p_a\log q_a,
\qquad p=\operatorname{softmax}(z_k+\delta z).
\]
因为 $p=\pi_k=q$ 时 KL 取全局最小值，$\nabla_zK(z_k)=0$。进一步使用
\[
\frac{\partial p_a}{\partial z_b}
=p_a(\mathbf 1\{a=b\}-p_b)
\]
可得 Hessian
\begin{equation}
\label{eq:ch2-logit-fisher}
\nabla_z^2K(z_k)
=
F_k(s)
=
\operatorname{Diag}(\pi_k)-\pi_k\pi_k^{\mathsf T}.
\end{equation}
$F_k$ 是 categorical policy 的 Fisher 矩阵，半正定且零空间由所有常数 logit 平移张成。因此
\begin{equation}
\label{eq:ch2-kl-logit-expansion}
\KL(\pi_\theta\Vert\pi_k)
=
\frac12\delta z^{\mathsf T}F_k\delta z
+R_{\mathrm{KL}}(\delta z),
\qquad
|R_{\mathrm{KL}}(\delta z)|
\leq C_{\mathrm{KL}}\|\delta z\|^3
\end{equation}
在 logits 位于有界集合时成立。这个三阶余项来自 softmax 的三阶导数有界，而不是额外的二元 KL 假设。

\paragraph{第二步：reward 项的一阶展开。}
令
\[
\ell_k(\theta)
=
-\E_{s,a}[r_k(s,a)]
\]
表示最小化目标中的 reward 部分。假设 $\ell_k$ 在 $\theta_k$ 邻域内二阶可微，则
\begin{equation}
\label{eq:ch2-reward-taylor}
\ell_k(\theta_k+\Delta\theta)
=
\ell_k(\theta_k)
-b_k^{\mathsf T}\Delta\theta
+R_r(\Delta\theta),
\qquad
b_k=\E_{s,a}\!\left[
r_k(s,a)\nabla_\theta\log\pi_{\theta_k}(a\mid s)
\right],
\end{equation}
并且 $|R_r(\Delta\theta)|\leq C_r\|\Delta\theta\|^2$。这里 $b_k$ 是固定 outer target 在当前学生处的 policy-gradient 方向；如果使用 PMD-Part 或 BPO 的回归 loss，则 $b_k$ 被相应的 regression score 替代，但 KL 曲率项相同。

\paragraph{第三步：从 logits Jacobian 得到参数空间曲率。}
令
\[
J_k(s)=\left.\frac{\partial z_\theta(s)}{\partial\theta}\right|_{\theta=\theta_k}.
\]
由
\[
\delta z
=
J_k(s)\Delta\theta+R_z(\Delta\theta),
\qquad
\|R_z(\Delta\theta)\|\leq C_z\|\Delta\theta\|^2,
\]
把式 \eqref{eq:ch2-kl-logit-expansion} 代入式 \eqref{eq:ch2-opd-inner-objective}，得到
\begin{equation}
\label{eq:ch2-kl-parameter-expansion}
\E_s\KL(\pi_\theta\Vert\pi_k)
=
\frac12\Delta\theta^{\mathsf T}A_k\Delta\theta
+R_{\mathrm{par}}(\Delta\theta),
\qquad
A_k=\E_s[J_k(s)^{\mathsf T}F_k(s)J_k(s)],
\end{equation}
其中 $|R_{\mathrm{par}}(\Delta\theta)|\leq C\|\Delta\theta\|^3$。因此式 \eqref{eq:ch2-local-quadratic-derived} 不是直接假设，而是式 \eqref{eq:ch2-kl-parameter-expansion} 与式 \eqref{eq:ch2-reward-taylor} 合并后，在 $\Delta\theta$ 足够小时的二阶主项：
\begin{equation}
\label{eq:ch2-local-quadratic-derived}
\mathcal J_k(\theta_k+\Delta\theta)
=
\mathcal J_k(\theta_k)
+\frac{1}{2\eta}\Delta\theta^{\mathsf T}A_k\Delta\theta
-b_k^{\mathsf T}\Delta\theta
+R_k(\Delta\theta).
\end{equation}

\begin{proposition}[局部曲率的正半定性与不可辨识方向]
\label{prop:ch2-local-curvature}
$A_k$ 对称半正定。若某个参数方向 $v$ 满足 $J_k(s)v$ 对所有 replay prefix 都是常数 logit 平移，则 $A_kv=0$；该方向不改变策略分布，因而 KL 不能为它提供曲率。若在商空间上 $A_k\succeq mI$，则镜像 inner 目标在该商空间局部强凸。
\end{proposition}

\begin{proof}
对任意 $v$，
\[
v^{\mathsf T}A_kv
=
\E_s[(J_k(s)v)^{\mathsf T}F_k(s)(J_k(s)v)]\geq0
\]
由 $F_k\succeq0$ 成立。$F_k$ 的零空间正是常数向量方向，所以 $J_k(s)v$ 为常数时对应二次型为零。去除这些不改变策略的参数方向后，若最小特征值至少为 $m>0$，则二次项给出局部强凸性。
\end{proof}

\subsection*{有限 inner 步的收敛条件}

定义局部主模型
\[
Q_k(\Delta\theta)
=
\frac{1}{2\eta}\Delta\theta^{\mathsf T}A_k\Delta\theta
-b_k^{\mathsf T}\Delta\theta.
\]
若 $b_k$ 属于 $A_k$ 的值域，则其最小范数解为
\[
\Delta\theta_k^\star=\eta A_k^\dagger b_k.
\]
使用普通梯度下降、步长 $\gamma>0$：
\begin{equation}
\label{eq:ch2-inner-gd-derived}
\Delta\theta_{m+1}
=
\Delta\theta_m-\gamma\nabla Q_k(\Delta\theta_m)
=
(I-\gamma A_k/\eta)\Delta\theta_m+\gamma b_k.
\end{equation}

\begin{theorem}[由 KL 曲率推出的 inner 线性收敛]
\label{thm:ch2-local-gd-derived}
设 $A_k$ 在商空间上的特征值属于
$[m_k,M_k]$，且 $0<m_k\leq M_k$。若
\[
0<\gamma<\frac{2\eta}{M_k},
\]
则式 \eqref{eq:ch2-inner-gd-derived} 收敛到 $\Delta\theta_k^\star$，并满足
\[
\|\Delta\theta_m-\Delta\theta_k^\star\|_{A_k}
\leq
q_k^m\|\Delta\theta_0-\Delta\theta_k^\star\|_{A_k},
\qquad
q_k=\max_{\lambda\in[m_k,M_k]}
\left|1-\frac{\gamma\lambda}{\eta}\right|<1.
\]
特别地，取 $\gamma=2\eta/(m_k+M_k)$ 时，
\[
q_k=\frac{M_k-m_k}{M_k+m_k}
=\frac{\kappa_k-1}{\kappa_k+1},
\qquad
\kappa_k=M_k/m_k.
\]
\end{theorem}

\begin{proof}
令 $e_m=\Delta\theta_m-\Delta\theta_k^\star$。由式 \eqref{eq:ch2-inner-gd-derived}，
\[
e_{m+1}=(I-\gamma A_k/\eta)e_m.
\]
在 $A_k$ 的特征方向 $\lambda$ 上，误差乘以 $1-\gamma\lambda/\eta$。给定步长条件，所有这些因子的绝对值都小于 $1$，因此得到 $q_k^m$ 的界。令两个端点方向的绝对收缩率相等即可得到最优固定步长和最后的条件数表达式。
\end{proof}

真实 inner 目标还包含余项 $R_k$。若其梯度在半径 $r$ 的邻域满足
\[
\|\nabla R_k(\Delta\theta)\|
\leq L_k\|\Delta\theta\|^2,
\]
并且 $r$ 足够小使得 $L_k r\leq m_k/(2\eta)$，则真实 Hessian 的特征值仍落在
\[
[m_k/(2\eta),\,M_k/\eta+L_kr]
\]
内。于是标准强凸光滑函数的梯度下降结论仍成立，只是收缩率由上述扩大后的光谱区间决定。这给出了“inner 需要做到局部收敛”的明确条件：步长必须由 KL/Fisher 曲率控制，且迭代必须停留在二阶展开有效的信赖域内。

\subsection*{Natural Gradient 的来源与作用}

普通梯度在参数坐标中使用 $A_k$ 的谱，病态方向的收敛速度由 $\kappa_k$ 决定。Natural Gradient 以 Fisher 度量约束一次更新：
\begin{equation}
\label{eq:ch2-natural-gradient-subproblem}
\min_{\Delta\theta}\ 
-b_k^{\mathsf T}\Delta\theta
\quad\text{s.t.}\quad
\frac12\Delta\theta^{\mathsf T}A_k\Delta\theta\leq\frac{\eta}{2}.
\end{equation}
拉格朗日一阶条件给出
\[
\Delta\theta_{\mathrm{NG}}
=
c\,A_k^\dagger b_k,
\]
其中常数 $c$ 由 trust-region 半径确定。若直接取 $\Delta\theta=\eta A_k^\dagger b_k$，它正是局部二次模型的最小点；实际算法用 Fisher 近似 $\widehat A_k$、阻尼 $\lambda I$ 和有限 inner 步，因此得到的是近似自然梯度，而不是理论上无条件的一步收敛。

\begin{proposition}[Natural Gradient 对参数重参数化的不变性]
\label{prop:ch2-ng-invariance}
设 $\vartheta=\phi(\theta)$ 是局部可逆的光滑重参数化，且 Fisher 矩阵按张量变换规则变换。则更新方向
\[
\Fisher_\theta^\dagger\nabla_\theta\mathcal J_k
\]
与在 $\vartheta$ 坐标中计算后映射回 $\theta$ 的方向一致；普通梯度 $\nabla_\theta\mathcal J_k$ 一般不具备该性质。
\end{proposition}

\begin{proof}
令 $B=\partial\vartheta/\partial\theta$。梯度变换为 $\nabla_\theta\mathcal J=B^{\mathsf T}\nabla_\vartheta\mathcal J$，Fisher 变换为 $\Fisher_\theta=B^{\mathsf T}\Fisher_\vartheta B$。在可逆子空间上代入并使用伪逆的相应变换，得到
\[
\Fisher_\theta^\dagger\nabla_\theta\mathcal J
=
B^{-1}\Fisher_\vartheta^\dagger\nabla_\vartheta\mathcal J.
\]
普通梯度只含 $B^{\mathsf T}$，因此会随坐标尺度改变。
\end{proof}

这解释了 MDPO 和 EffOPD 中的两个实现判断。只在 $\theta_k$ 处做一次梯度步时，KL 的一阶梯度为零，trust-region 曲率尚未进入更新；固定 $\pi_k$ 并执行多次 inner 更新，才能逼近式 \eqref{eq:ch2-opd-inner-objective}。Natural Gradient 进一步把步长按照 $A_k$ 的几何尺度归一化，但它仍然需要信赖域、阻尼和有限 inner 预算来处理模型非线性与 Fisher 估计误差。


\paragraph{几何中间分布与算术 proximal teacher 的区别。} 式 \eqref{eq:ch2-opd-geometric-bridge} 给出的 $\pi_k^{1-\eta}T^\eta/Z_k$ 是 KL 镜像子问题的精确解，因此在 log-probability 空间插值。TOP-D 中的算术 teacher
\[
\widetilde T_\eta=(1-\eta)\pi_k+\eta T
\]
是另一个用于控制 log-ratio 方差的 teacher operator。它同样产生 student 与 teacher 之间的中间分布，但不等价于 KL mirror 解；前者改变 target teacher，后者改变 proximal 优化问题的解。

